\subsection{Uniqueness by concentration of mass}
\label{sec:relative-prescribed-measure-uniqueness}
Existence for every matching non-pluripolar measure lets us concentrate an auxiliary solution's mass on one strict comparison set. The mixed inequality then increases the density on that set. Full-mass comparison is incompatible with this increase unless the two original potentials differ by a constant almost everywhere for their common measure; domination gives pointwise equality up to that constant.

\begin{theorem}\label{thm:relative-prescribed-measure-uniqueness}
Keep the big background $\theta$, the model reference $\phi$ and $m=M(\phi)>0$. If $u,v\in\mathcal E(X,\theta,\phi)$ satisfy $\mu_u=\mu_v$, then there is $C\in\mathbb R$ such that $u=v+C$ everywhere. Thus each non-pluripolar measure of mass $m$ has exactly one solution in this relative full-mass class with supremum zero.
\end{theorem}
\begin{proof}
Put $\mu=\mu_u=\mu_v$. This is a finite non-pluripolar measure of mass $m$. Both potentials are finite $\mu$-almost everywhere. Suppose their difference is not constant $\mu$-almost everywhere. There are two distinct essential values of that finite difference; between them choose a level which is not an atom of its distribution. Only countably many levels are atoms. After adding this level to $v$, we have
\[
 U=\{u<v\},\qquad0<\mu(U)<m,
                     \qquad\mu(\{u=v\})=0.
\]
Constant shifts preserve the relative class and the measures. Set $q=(m/\mu(U))^{1/n}>1$. By \cref{thm:nonpluripolar-relative-prescribed-measure-existence}, choose $h\in\mathcal E(X,\theta,\phi)$ with
\[
 \mu_h=q^n\mathbf1_U\mu.
\]
Its prescribed measure is non-pluripolar and has mass $m$.

For $0<t<1$ put $w_t=(1-t)v+th$. Convexity, \cref{cor:relative-energy-classes-convex}, gives $w_t\in\mathcal E(X,\theta,\phi)$. Apply \cref{cor:nonpluripolar-common-measure-convex} to $\theta_v$ and $\theta_h$, with common measure $\mu$ and densities $1,q^n\mathbf1_U$. It gives
\[
 \mu_{w_t}\geq\big((1-t)+tq\mathbf1_U\big)^n\mu.
\]
Let $U_t=\{u<w_t\}$ and write
\[
 d_t=(1-t)^n,\qquad b_t=(1-t+tq)^n.
\]
Full-mass comparison, \cref{prop:relative-mixed-comparison}, gives
\[
 \mu(U_t)\geq\mu_{w_t}(U_t)
       \geq d_t\mu(U_t)+(b_t-d_t)\mu(U_t\cap U).
\]
Since $1-d_t>0$, we obtain
\begin{equation}\label{eq:relative-uniqueness-concentration}
 \frac{b_t-d_t}{1-d_t}\,\mu(U_t\cap U)\leq\mu(U_t).
\end{equation}
The three potentials $u,v,h$ are finite $\mu$-almost everywhere, and the level $u=v$ is null. Consequently $\mathbf1_{U_t}\to\mathbf1_U$ $\mu$-almost everywhere as $t\downarrow0$. Dominated convergence gives
\[
 \mu(U_t\cap U),\ \mu(U_t)\longrightarrow\mu(U).
\]
Also $(b_t-d_t)/(1-d_t)\to q$, by differentiating the two polynomials at zero. Passing to the limit in \eqref{eq:relative-uniqueness-concentration} gives $q\mu(U)\leq\mu(U)$, contradicting $q>1$ and $\mu(U)>0$.

Thus for the original pair $u=v+C$ $\mu$-almost everywhere, with a finite $C$. Both strict inequality sets for $u,v+C$ are $\mu$-null. Apply \cref{thm:relative-full-mass-domination} in each direction to obtain $u=v+C$ everywhere, including at common poles. Supremum-zero normalization forces $C=0$; the preceding existence theorem supplies the solution. The same proof applies in dimension one.
\end{proof}
\index{Monge--Amp\`ere equation!relative prescribed-measure uniqueness}
